\section{Conclusion}

This paper advances a deliberately narrow claim. It does not propose a general governance platform for FDO-based agent systems. It proposes a minimal verifiable profile for single-operation accountability, together with a profile-aware validator and an implementation-grounded artifact package. In the current implementation, that claim is supported by the profile specification, the JSON Schema, 2 valid examples, 5 invalid examples, the validator and command-line path, a runnable single-path demo, and a frozen review package with public archive identifiers withheld for anonymous review.

This work reduces operation accountability to a form that can be specified, checked, reproduced, and archived before broader system claims are made. The paper therefore argues for a small engineering object with explicit links among \texttt{operation}, \texttt{policy}, \texttt{provenance}, \texttt{evidence}, and \texttt{validation}, rather than for a broad theory of accountability across all agent-system settings.

What the paper does not claim is equally important. It does not claim broad cross-framework validation, industrial deployment, complete workflow coverage, or a full cryptographic trust infrastructure. The current implementation evidence supports a narrower conclusion: one operation accountability statement can be represented as a minimal verifiable profile with a concrete validation path and a reproducible artifact package.

Near-term follow-on work should remain on the same mainline: add a small number of controlled examples, strengthen validator tests, complete the remaining citation and packaging cleanup, and keep the profile boundary stable while the evidence base is thickened. The next step is not to inflate the scope into a broad governance framework, but to reinforce the same minimal profile, validator, and artifact package with incrementally stronger implementation-grounded evidence.
